\section{Overview of our Techniques}

{\bf Proof Outline of Theorem~\ref{thm:supervised-distill-1}.} The proof proceeds by contradiction: suppose there is some $h \in \mc{F}$ and $D^{\sf sup}_{\tn{syn}}$ s.t. $L^{\tn{err}}_{\tn{mse}}(D^{\sf sup}_{\tn{train}}, D^{\sf sup}_{\tn{syn}}, h) > \Delta$. Letting $g$ be a random sample from  $G$, we show  using algebraic manipulations and properties of the Gaussian distribution that  $L^{\tn{err}}_{\tn{mse}}(D^{\sf sup}_{\tn{train}}, D^{\sf sup}_{\tn{syn}}, g)$ is a degree-$4$ polynomial $\upsilon^2$ in Gaussian variables, such that $\E[\upsilon^2] > \Delta/O(d^2)$. This can be used along with the Carbery-Wright anti-concentration bound for Gaussian polynomials to show that $\Pr[\upsilon^2 > \Delta/O(d^2)] \geq 1/3$. Since $\{g_j\}_{j=1}^k$ are iid samples from $G$, using Chernoff Bound, the probability that there exists $\Omega(k)$ many $j \in [k]$ satisfying  $L^{\tn{err}}_{\tn{mse}}(D^{\sf sup}_{\tn{train}}, D^{\sf sup}_{\tn{syn}}, g_j) > \Delta/O(d^2)$ is at least $1 - \tn{exp}^{-\Omega(k)}$. To complete the argument, we require a union bound of this error probability over all $h \in \mc{F}$ and $D^{\sf sup}_{\tn{syn}}$, which we do by constructing finegrained nets as follows.\\
{\it Net over all $h \in \mc{F}$}: Since $h \in \mc{F}$ is given by some $\br \in \R^{d+1}$ where $1 \leq \|\br\|_2 \leq 2$, one can take a $\eps$-net w.r.t. to Euclidean metric over all such vectors -- whose size is at most $(1/\eps)^{O(d\log d)}$ (see Appendix \ref{sec:nets}) -- for a small enough $\eps$ so that $L^{\tn{err}}_{\tn{mse}}$ is essentially unaffected when $h$ is replaced by $\hat{h}$ (or vice versa) corresponding to the nearest vector in the net. For our purpose $\eps$ can be $\Omega(\Delta)$. 
 \\
{\it Net over all $D^{\sf sup}_{\tn{syn}}$}: For this we first observe 
% using pseudo-dimension based generalization error bounds 
that $D^{\sf sup}_{\tn{syn}}$ can always replaced by a subset of size
$s = d$ using Lemma~\ref{lem:subsetreplace} so that
% $s = O((d^5/\Delta^2)(\log(d/\Delta))$ so that
$L^{\tn{err}}_{\tn{mse}}$ remains lower bounded by $\Delta/O(d^2)$. Thus, it suffices to consider the net $(\mc{N})^{s}$ over all $s$-sized datasets where $\mc{N}$ is a  $\eps'$-net over the set of all possible datapoints,  so that $L^{\tn{err}}_{\tn{mse}}$ is approximately preserved by the nearest dataset in $(\mc{N})^{s}$. \\
% The final bound is obtained by setting the granularity of the two nets appropriately, and taking $k$ (the number of sampled regressors) to be large enough so that a union bound over the error can be taken.
The size of $(\mc{N})^{s}$ is $((B+b)/\Delta)^{O(sd)}$, and the product of the sizes of the nets constructed above is exponential in $O\left(d^2\log(d(B+b)/\Delta)\right)$. Thus, to obtain the statement of the theorem via a union bound, the number of sampled regressors required is 
$k = O\left(d^2\log(d(B+b)/\Delta)\log(1/\delta)\right)$. 

{\bf Proof Outline of Theorem~\ref{thm:lbnumber}.} The proof of the lower bound of $\Omega(d^2)$ on the sampled linear regressors relies on the fact that the set of symmetric $q\times q$ matrices, where $q=d+1$, is a $q(q+1)/2$-dimensional linear space.  Thus, for a choice of less than $q(q+1)/2$ homogeneous linear regressors, there is a non-zero symmetric matrix $\mb{A}$ for which $\langle \bv\bv^{\sf T}, \mb{A}\rangle = 0$ where $\bv$ is any one of the chosen regression vectors. Here, $\mb{A}$ can be scaled so that its operator norm is exactly $1/2$. Thus, one can choose (up to appropriate scaling) $D^{\sf sup}_{\tn{train}}$ so that $\E_{\bz \in D^{\sf sup}_{\tn{train}}}[\bx\bx^{\sf T}] = \mb{I}$ which is independent of the chosen regressors. $D^{\sf sup}_{\tn{syn}}$ is then chosen so that $\E_{\bz \in D^{\sf sup}_{\tn{syn}}}[\bx\bx^{\sf T}] = \mb{I} +\mb{A}$ which is psd. It can be seen that each chosen regressor has the same MSE loss on $D^{\sf sup}_{\tn{train}}$ and $D^{\sf sup}_{\tn{syn}}$, while due to $\mb{A} \neq \mb{0}$, there is a regressor which has different losses on $D^{\sf sup}_{\tn{train}}$ and $D^{\sf sup}_{\tn{syn}}$.\\


{\bf Proof Outline of Theorem~\ref{thm:offlineRL-distill-1}.}

The main complication as compared to the supervised regression case is the presence of the $\max$ term in the Bellman loss, which we circumvent using a conditioning argument. Specifically, consider some action-value predictor $f\in \mc{Q}_0$ given by $f(s, a) := \br^{\sf T}\phi(s,a)$ such that $\hat{L}^{\tn{err}}_{\tn{Bell}}(D^{\sf orl}_{\tn{train}}, D^{\sf orl}_{\tn{syn}}, f, 1) > \eps$, for some $D^{\sf orl}_{\tn{syn}}$. Assume for ease of exposition that $\|\br\|_2 = 1$. Now, a $\hat{f} \sim H$ corresponds to a vector $\hat{\br} = \alpha\br + \mb{J}\bu$ where $\alpha \sim N(0,1)$, $\bu \sim N(0,1)^{d-1}$ and $\mb{J}$ is a $d\times (d-1)$ matrix with columns being a completion of $\br$ to an orthonormal basis. It is easy to see that with probability at least $\tn{exp}(-{O(d\log d)})$, $\|\bu\|_2 \ll 1$ and therefore $\|\mb{J}\bu\|_2 \ll 1$. Conditioning on this event and the positivity of $\alpha$ allows us to use $\max_{a'}\hat{f}(\phi(s',a')) \approx \alpha\max_{a'}f(\phi(s',a'))$. This, along with a conditioning on $\lambda \approx 1$ and some algebraic manipulations yields that $\hat{L}^{\tn{err}}_{\tn{Bell}}(D^{\sf orl}_{\tn{train}}, D^{\sf orl}_{\tn{syn}}, \hat{f}) \approx \alpha^4 \Delta$, which directly yields a probabilistic lower bound of $\Omega(\Delta)$ on $\hat{L}^{\tn{err}}_{\tn{Bell}}(D^{\sf orl}_{\tn{train}}, D^{\sf orl}_{\tn{syn}}, \hat{f})$. 
The rest of the net based arguments are analogous to the supervised case but with notable differences. 
{In particular, we observe that $D_{\tn{syn}}^{\sf orl}$ can always be replaced by a subset of size $s= O(d/\Delta^2 \log{(1/\Delta)})$ so that $L_{\tn{Bell}}^{\tn{err}}$ remains lower bounded by $O(\Delta)$.}
In addition, we use a more broadly applicable generalization error bound for the the Bellman loss using its pseudo-dimension and the Lipschitzness of $\phi$. 
However, due to the $\tn{exp}(-{O(d\log d)})$ of the conditioning, success probability for a single sample $(\hat{f}, \lambda)$ is also $\tn{exp}(-{O(d\log d)})$ and therefore the number of predictors to be sampled is $\tn{exp}({O(d\log d)})$.

{\bf Proof outline of Theorem \ref{thm:offline-rl-decomposable}.} We observe that when $\phi$ is decomposable, for a $Q$-value predictor $f(s,a) := \bv^{\sf T}\phi(s,a)$, we have that $\max_{a'\in \mc{A}}f(s',a') = \bv^{\sf T}\phi_1(s') + \max_{a'\in \mc{A}}\bv^{\sf T}\phi_2(a')$. Here $\alpha := \max_{a'\in \mc{A}}\bv^{\sf T}\phi_2(a')$ depends only on $\bv$ and is independent of the dataset. This allows for cancellations of terms between $D^{\sf orl}_{\tn{train}}$ and $D^{\sf orl}_{\tn{syn}}$ in $\hat{L}^{\tn{err}}_{\tn{Bell}}$, and using the constraint in \eqref{eqn:meancondition} we obtain a linear regression formulation in the embedding space $\R^d$. Thus, one can essentially follow the proof of Theorem \ref{thm:supervised-distill-1}. Further, it is easy to see that if $\phi_1$ and $\phi_2$ are linear then the constraint in \eqref{eqn:meancondition} is linear in the points of $D^{\sf orl}_{\tn{syn}}$ resulting in a convex optimization.